\sloppy
\emergencystretch=5em
\hbadness=10000
\vbadness=10000
\tolerance=2000

\providecolor{codebg}{RGB}{248,248,248}
\providecolor{codeframe}{RGB}{200,200,200}
\providecolor{keyword}{RGB}{0,0,180}
\providecolor{string}{RGB}{163,21,21}
\providecolor{comment}{RGB}{0,128,0}

\lstset{
	backgroundcolor=\color{codebg},
	frame=single,
	rulecolor=\color{codeframe},
	language=Python,
	basicstyle=\ttfamily\scriptsize,
	keywordstyle=\color{keyword}\bfseries,
	stringstyle=\color{string},
	commentstyle=\color{comment}\itshape,
	showstringspaces=false,
	breaklines=true,
	breakatwhitespace=true,
	numbers=left,
	numberstyle=\tiny\color{gray},
	numbersep=8pt,
	tabsize=4,
	literate=
	{ف}{{{\bfseries\itshape ف}}}1
	{ی}{{{\bfseries\itshape ی}}}1
}

\section{نمای کلی و هدف ماژول}

ماژول \lr{\texttt{pulse\_tracking\_trace.py}} یک اسکریپت \emph{ممیزی و ردیابی سرتاسری} (\lr{end-to-end audit}) است که نقش یک \lr{forensic tracer} را در فریم‌ورک شبیه‌سازی \lr{QKD} ایفا می‌کند. این ماژول به‌طور خاص برای پاسخ به پرسش بنیادین زیر طراحی شده است: \emph{وقتی هزاران پالس نوری از طریق خط لوله چهارمرحله‌ای شبیه‌سازی (\lr{dispatch}، \lr{worker}، \lr{receipt}، \lr{row}) عبور می‌کنند، آیا هویت فیزیکی هر پالس — شامل بیت ارسالی \lr{Alice}، پایه انتخابی، پایه گیرنده \lr{Bob} و حالت کوانتومی — به‌طور سالم و بدون تغییر از مرحله‌ای به مرحله بعد منتقل می‌شود، یا این که در طول مسیر به‌دلیل \lr{reordering} موازی، \lr{race condition}، یا خطاهای پیاده‌سازی دچار فساد (\lr{corruption}) می‌گردد؟} این پرسش در شبیه‌سازی‌های موازی که در آن \lr{ProcessPoolExecutor} با چند کارگر روی هم‌زمان (\lr{concurrency}) اجرا می‌شود، حیاتی است؛ زیرا ترتیب خروجی (\lr{arrival order}) با ترتیب ورودی (\lr{input order}) لزوماً یکسان نیست و این ناهماهنگی می‌تواند به‌طور پنهانی منجر به جابجایی هویت پالس‌ها شود.

هدف اصلی این ماژول، تولید یک \emph{رد Tracks قابل ممیزی} (\lr{auditable trace}) است که برای ۱۰۰ پالس تصادفی — از جمعیت ۱۰۰۰ پالس — نشان می‌دهد که دقیقاً چه بر سر هر پالس در هر یک از چهار مرحله می‌آید. چهار مرحله عبارت‌اند از: مرحله \lr{S1} (\lr{Pre-Dispatch Snapshot}) که هویت اصلی پالس پیش از ارسال به کارگران را ثبت می‌کند؛ مرحله \lr{S2} (\lr{Receipt Snapshot}) که هویت پالس پس از دریافت در کارگر را نشان می‌دهد؛ مرحله \lr{S3} (\lr{Worker Output}) که نتیجه اندازه‌گیری کارگر را در بر دارد؛ و مرحله \lr{S4} (\lr{Final Row}) که ردیف نهایی \lr{CSV} را نمایش می‌دهد. اگر هر چهار مرحله بر سر تمام \lr{identity fields} توافق کنند، یعنی هویت پالس به‌طور کامل حفظ شده است؛ در غیر این صورت، خطای فساد رخ داده است.

ساختار کلی این ماژول شامل چندین بخش است. بخش اول، توابع قالب‌بندی (\lr{formatting helpers}) شامل \lr{\texttt{\_format\_table}}، \lr{\texttt{\_short\_identity}}، \lr{\texttt{\_identity\_match}} و \lr{\texttt{\_measurement\_match}} است که مسئول تولید خروجی متنی خوانا هستند. بخش دوم، تابع اصلی \lr{\texttt{run\_trace}} است که جریان کامل آزمون را هماهنگ می‌کند: ساخت جمعیت پالس، انتخاب پالس‌های ردیابی‌شده، اجرای \lr{dispatch}، محاسبه اثبات \lr{reordering}، تولید جدول ممیزی چهارمرحله‌ای، و محاسبه آمار سطح-مرحله‌ای. بخش سوم، تابع \lr{\texttt{main}} و \lr{CLI} است که از \lr{\texttt{argparse}} برای پیکربندی پارامترها (تعداد کارگران، اندازه جمعیت، تعداد پالس‌های ردیابی‌شده، بذر تصادفی، تعداد ردیف‌های نمایش‌داده‌شده، فایل خروجی) استفاده می‌کند.

ویژگی برجسته این ماژول، الگوی \lr{reuse-by-import} آن است؛ یعنی به‌جای بازنویسی ماشین‌آلات آزمون، مستقیماً از توابع کمکی در ماژول \lr{\texttt{tests/test\_main\_optimized\_pulse\_tracking.py}} استفاده می‌کند. این توابع شامل \lr{\texttt{\_make\_pulse\_population}} برای ساخت جمعیت، \lr{\texttt{\_select\_tracked\_pids}} برای انتخاب تصادفی پالس‌های ردیابی‌شده با بذر قابل بازتولید، \lr{\texttt{\_run\_tracking\_dispatch}} برای اجرای \lr{dispatch} با ردیابی، \lr{\texttt{\_expected\_measurement}} برای پیش‌بینی نتیجه اندازه‌گیری بر اساس هویت پالس، و \lr{\texttt{\_tracking\_slow\_worker}} که یک کارگر با تأخیر دوجمله‌ای (\lr{bimodal delay}) است. این الگو، اصل \lr{DRY} (\lr{Don't Repeat Yourself}) را رعایت می‌کند و تضمین می‌نماید که اسکریپت ممیزی و آزمون واحد، از همان ماشین‌آلات استفاده کنند و در نتیجه، نتایج آن‌ها همیشه هم‌خوان باشند.

در سطح عملیاتی، این ماژول به‌عنوان یک \lr{diagnostic tool} عمل می‌کند؛ یعنی در صورت بروز مشکلات مربوط به حفظ هویت پالس (مثلاً پس از \lr{refactor} موتور \lr{dispatch} یا تغییر مدل کارگر)، توسعه‌دهنده می‌تواند این اسکریپت را اجرا کند و به‌سرعت علت فساد را شناسایی کند. خروجی آن به دو شکل است: چاپ به \lr{stdout} برای مشاهده زنده، و ذخیره به فایل \lr{\texttt{pulse\_tracking\_trace.txt}} در مسیر \lr{\texttt{/home/z/my-project/download/}} برای بازبینی بعدی. این الگو در مهندسی نرم‌افزار علمی به‌عنوان \lr{shift-left forensics} شناخته می‌شود و به توسعه‌دهنده اجازه می‌دهد خطاهای ظریف را پیش از انتشار نتایج شناسایی کند.

نکته مهم دیگر آن است که این ماژول \emph{اثبات \lr{reordering}} (\lr{reordering proof}) را به‌صورت صریح و قابل مشاهده ارائه می‌دهد. در شبیه‌سازی موازی، خروجی کارگران به ترتیب ورودی نیست؛ بلکه ترتیب ورودی به ترتیب اتمام کارگران بستگی دارد. این ماژول با محاسبه تعداد \lr{pairwise inversions} (یعنی زوج‌های پالسی که ترتیب ورودی‌شان با ترتیب خروجی متفاوت است)، نشان می‌دهد که \lr{reordering} رخ داده اما هیچ هویتی فاسد نشده است. این تمایز بنیادینی است: \lr{reordering} یک ویژگی ذاتی و قابل‌قبول شبیه‌سازی موازی است، در حالی که \lr{corruption} یک باگ است که باید برطرف شود.

\section{مبانی تئوری و فیزیکی}

\subsection{پروتکل BB84 و مفهوم هویت پالس}

پروتکل \lr{BB84} که توسط \lr{Bennett} و \lr{Brassard} در سال ۱۹۸۴ پیشنهاد شد، ساده‌ترین و پرکاربردترین پروتکل \lr{QKD} است. در این پروتکل، فرستنده (\lr{Alice}) یک بیت کلاسیک را با انتخاب تصادفی یکی از دو پایه (\lr{basis}) — پایه مستطیلی (\lr{rectilinear}) $\{|0\rangle, |1\rangle\}$ یا پایه قطری (\lr{diagonal}) $\{|+\rangle, |-\rangle\}$ — به گیرنده (\lr{Bob}) ارسال می‌کند. \lr{Bob} نیز به‌طور تصادفی یکی از این دو پایه را برای اندازه‌گیری انتخاب می‌کند. در صورتی که پایه‌های \lr{Alice} و \lr{Bob} یکسان باشند، نتیجه اندازه‌گیری \lr{Bob} با بیت ارسالی \lr{Alice} هم‌خوان است (به جز خطاهای ناشی از نویز کانال و آشکارساز)؛ در غیر این صورت، نتیجه تصادفی است.

هویت هر پالس در شبیه‌سازی \lr{BB84} با چندین فیلد کلیدی مشخص می‌شود که در مجموع \emph{identity fields} نامیده می‌شوند. این فیلدها در ثابت \lr{\texttt{IDENTITY\_FIELDS}} در ماژول آزمون تعریف شده‌اند و شامل موارد زیر هستند: \lr{\texttt{alice\_bit}} (بیت ارسالی \lr{Alice}، صفر یا یک)، \lr{\texttt{alice\_basis}} (پایه انتخابی \lr{Alice}، مستطیلی یا قطری)، \lr{\texttt{bob\_basis}} (پایه انتخابی \lr{Bob}، مستطیلی یا قطری)، \lr{\texttt{state\_label}} (برچسب حالت کوانتومی مانند \lr{|0⟩}، \lr{|1⟩}، \lr{|+⟩}، \lr{|−⟩})، \lr{\texttt{mu}} (شدت میانگین پالس) و \lr{\texttt{intensity\_label}} (برچسب شدت مانند \lr{signal}، \lr{decoy} یا \lr{vacuum}). حفظ یکپارچگی این فیلدها در طول خط لوله شبیه‌سازی، شرط لازم برای صحت نتایج است.

\subsection{پیش‌بینی نتیجه اندازه‌گیری و قانون Born}

نتیجه اندازه‌گیری یک پالس \lr{BB84} توسط \lr{Bob} را می‌توان بر اساس هویت پالس پیش‌بینی کرد. این پیش‌بینی بر اساس قانون \lr{Born} و قواعد پروتکل \lr{BB84} انجام می‌شود. در صورتی که پایه‌های \lr{Alice} و \lr{Bob} یکسان باشند (\lr{basis\_match = True}) و پالس به‌طور موفقیت‌آمیزی شناسایی شود (\lr{detected = True})، بیت اندازه‌گیری‌شده با بیت ارسالی هم‌خوان است (\lr{measured\_bit = alice\_bit}). در صورتی که پایه‌ها متفاوت باشند، نتیجه اندازه‌گیری کاملاً تصادفی است و احتمال آن $0.5$ برای هر بیت است. در شبیه‌سازی قطعی با بذر مشخص، این تصادفی بودن از یک \lr{PRNG} قابل بازتولید به دست می‌آید.

پیش‌بینی نتیجه اندازه‌گیری در تابع \lr{\texttt{\_expected\_measurement}} پیاده‌سازی شده است. این تابع یک \lr{dict} شامل سه کلید \lr{\texttt{basis\_match}}، \lr{\texttt{measured\_bit}} و \lr{\texttt{detected}} برمی‌گرداند که مقادیر پیش‌بینی‌شده را در بر دارد. مقایسه این مقادیر با مقادیر گزارش‌شده توسط کارگر (\lr{S3})، امکان تشخیص خطاهای پیاده‌سازی را فراهم می‌سازد. به‌عنوان مثال، اگر کارگر \lr{\texttt{basis\_match = True}} و \lr{\texttt{detected = True}} را گزارش دهد اما \lr{\texttt{measured\_bit}} با \lr{\texttt{alice\_bit}} متفاوت باشد، این یک ناسازگاری است که نشان‌دهنده باگ در منطق اندازه‌گیری کارگر است.

\subsection{تأخیر دوجمله‌ای و شبیه‌سازی کارگر واقعی}

یکی از چالش‌های کلیدی در آزمون شبیه‌سازی موازی، ایجاد شرایط واقعی برای آشکار کردن \lr{race condition} و \lr{reordering} است. اگر همه کارگران با تأخیر یکسان اجرا شوند، خروجی به‌طور قطعی به ترتیب ورودی است و \lr{reordering} رخ نمی‌دهد. برای غلبه بر این مشکل، تابع \lr{\texttt{\_tracking\_slow\_worker}} با تأخیر دوجمله‌ای (\lr{bimodal delay}) پیاده‌سازی شده است:

\begin{equation}
	T_{\text{worker}} = \begin{cases} 100 \text{ ms} & \text{با احتمال } 0.10 \\ U(0, 5) \text{ ms} & \text{با احتمال } 0.90 \end{cases}
\end{equation}

در این رابطه، $T_{\text{worker}}$ زمان اجرای کارگر است و $U(0, 5)$ توزیع یکنواخت در بازه $[0, 5]$ میلی‌ثانیه است. این مدل، دو رژیم را شبیه‌سازی می‌کند: $10\%$ کارگران \emph{کند} هستند (تأخیر ۱۰۰ میلی‌ثانیه که شبیه‌سازی بار \lr{CPU} سنگین یا \lr{I/O} است) و $90\%$ کارگران \emph{سریع} هستند (تأخیر تصادفی ۰ تا ۵ میلی‌ثانیه). این توزیع غیریکنواخت، احتمال \lr{reordering} را به‌طور قابل‌توجهی افزایش می‌دهد و امکان آشکارسازی باگ‌های مربوط به حفظ ترتیب را فراهم می‌سازد.

\subsection{ Kendall-Tau و معیار ناهماهنگی}

برای اندازه‌گیری کمی میزان \lr{reordering}، از معیار \lr{Kendall-Tau inversion count} استفاده می‌شود. این معیار، تعداد زوج‌های (\lr{i, j}) با $i < j$ را می‌شمارد که در آن ترتیب خروجی برعکس ترتیب ورودی است:

\begin{equation}
	\text{inv}(\pi) = \left| \{(i, j) : i < j \land \pi(i) > \pi(j)\} \right|
\end{equation}

در این رابطه، $\pi$ جایگشت خروجی نسبت به ورودی است. حداکثر تعداد وارونگی‌ها برای $n$ پالس برابر با $\binom{n}{2} = \frac{n(n-1)}{2}$ است که در صورت جایگشت معکوس کامل رخ می‌دهد. درصد وارونگی به‌صورت زیر تعریف می‌شود:

\begin{equation}
	\text{reorder\%} = \frac{\text{inv}(\pi)}{\binom{n}{2}} \times 100
\end{equation}

این معیار، یک اندازه عاری از مقیاس (\lr{scale-free}) برای سنجش شدت \lr{reordering} فراهم می‌سازد. در شبیه‌سازی با ۱۰۰ پالس ردیابی‌شده و چند کارگر با تأخیر دوجمله‌ای، مقدار معمول این معیار در محدوده $30\%$ تا $50\%$ است که نشان‌دهنده \lr{reordering} قابل‌توجه اما نه کامل است.

\subsection{مفهوم Identity-Bound-to-ID و Position-Independent Integrity}

مفهوم کلیدی که این ماژول آن را به‌اثبات می‌رساند، \emph{یکپارچگی وابسته به شناسه، نه به موقعیت} (\lr{identity bound to pulse\_id, not to row position}) است. در یک طراحی ساده اما خطاکار، ممکن است هویت پالس بر اساس موقعیت ردیف در خروجی استنتاج شود؛ یعنی فرض شود که ردیف $i$ در خروجی، مربوط به پالس $i$ در ورودی است. این فرض در حضور \lr{reordering} نادرست است و منجر به تخصیص اشتباه هویت به ردیف‌ها می‌شود.

طراحی صحیح، هویت را به \lr{\texttt{pulse\_id}} (یک شناسه یکتا که در زمان ساخت پالس تخصیص می‌یابد) می‌بندد. این شناسه به‌صورت یکتا در کل جمعیت است و در طول خط لوله بدون تغییر منتقل می‌شود. هر مرحله از خط لوله، این شناسه را در خروجی خود شامل می‌شود و امکان بازسازی هویت پالس را فراهم می‌سازد. این الگو در ادبیات مهندسی نرم‌افزار به‌عنوان \lr{correlation ID} یا \lr{trace ID} شناخته می‌شود و در سیستم‌های توزیع‌شده (\lr{distributed systems}) برای ردیابی درخواست‌ها به‌کررات به‌کار می‌رود.

\section{تحلیل معماری و ساختار داده‌ها}

\subsection{ساختار کلی ماژول و الگوی Script-with-Helpers}

این ماژول از الگوی \lr{script-with-helpers} پیروی می‌کند؛ یعنی یک تابع اصلی \lr{\texttt{run\_trace}} که جریان اجرا را هماهنگ می‌کند، و چندین تابع کمکی برای قالب‌بندی و مقایسه. این الگو برای اسکریپت‌های \lr{CLI} مناسب است که نیاز به خوانایی بالا و قابلیت آزمون دارند. تابع \lr{\texttt{run\_trace}} خروجی را به‌صورت یک رشته متنی بزرگ برمی‌گرداند (نه چاپ مستقیم)، که امکان نوشتن آن در \lr{stdout} یا فایل را فراهم می‌سازد. این جداسازی بین تولید خروجی و چاپ، الگوی \lr{separation of concerns} را رعایت می‌کند.

\subsection{ثابت IDENTITY\_FIELDS و فیلدهای هویتی}

ثابت \lr{\texttt{IDENTITY\_FIELDS}} که از ماژول آزمون وارد می‌شود، فهرست فیلدهایی است که برای تعیین هم‌خوانی هویت در چهار مرحله بررسی می‌شوند. این فیلدها در ساختار \lr{dict} هر پالس وجود دارند و شامل \lr{\texttt{alice\_bit}}، \lr{\texttt{alice\_basis}}، \lr{\texttt{bob\_basis}}، \lr{\texttt{state\_label}} و سایر فیلدهای هویتی هستند. استفاده از یک ثابت متمرکز برای این فیلدها، تضمین می‌کند که تمام بخش‌های ماژول از همان فهرست استفاده کنند و در صورت افزودن فیلد جدید، فقط یک نقطه نیاز به به‌روزرسانی دارد.

\begin{LTR}
	\begin{lstlisting}[language=Python]
		from tests.test_main_optimized_pulse_tracking import (
		IDENTITY_FIELDS,
		_make_pulse_population,
		_select_tracked_pids,
		_run_tracking_dispatch,
		_expected_measurement,
		_tracking_slow_worker,
		)
		from main_optimized_dispatch import dispatch_work_items
	\end{lstlisting}
\end{LTR}

این الگوی \lr{import} چندین نکته دارد. نخست، \lr{\texttt{IDENTITY\_FIELDS}} به‌عنوان یک ثابت وارد می‌شود تا در تمام ماژول قابل استفاده باشد. دوم، توابع کمکی با پیشوند \lr{\_} (مانند \lr{\texttt{\_make\_pulse\_population}}) وارد می‌شوند که در پایتون نشان‌دهنده \lr{private API} است؛ اما چون در همان پروژه هستیم، این واردسازی مجاز است. سوم، \lr{\texttt{dispatch\_work\_items}} از ماژول \lr{\texttt{main\_optimized\_dispatch}} وارد می‌شود که این نشان‌دهنده وابستگی مستقیم به موتور \lr{dispatch} اصلی فریم‌ورک است.

\subsection{تابع \_format\_table و قالب‌بندی متنی}

تابع \lr{\texttt{\_format\_table}} یک ابزار قالب‌بندی متنی است که لیستی از سطرها را به جدولی با ستون‌های هم‌تراز تبدیل می‌کند. این تابع دو پارامتر اصلی دارد: \lr{\texttt{headers}} (فهرست عناوین ستون‌ها) و \lr{\texttt{rows}} (لیست سطرها، که هر سطر یک لیست از مقادیر است). پارامتر سوم \lr{\texttt{col\_widths}} اختیاری است و در صورت عدم ارائه، به‌صورت خودکار بر اساس حداکثر طول مقدار در هر ستون محاسبه می‌شود.

\begin{LTR}
	\begin{lstlisting}[language=Python]
		def _format_table(headers: List[str], rows: List[List[Any]],
		col_widths: List[int] = None) -> str:
		"""Format a list of rows as an aligned text table."""
		if not rows:
		return "(no rows)"
		if col_widths is None:
		col_widths = []
		for i, h in enumerate(headers):
		w = max(len(str(h)), *(len(str(r[i])) for r in rows))
		col_widths.append(w)
		lines = []
		lines.append("  ".join(str(h).ljust(w) for h, w in zip(headers, col_widths)))
		lines.append("  ".join("-" * w for w in col_widths))
		for r in rows:
		lines.append("  ".join(str(c).ljust(w) for c, w in zip(r, col_widths)))
		return "\n".join(lines)
	\end{lstlisting}
\end{LTR}

در خط اول، شرط \lr{\texttt{not rows}} بررسی می‌شود و در صورت خالی بودن، رشته \lr{"(no rows)"} بازگردانده می‌شود. این الگو، الگوی \lr{graceful handling of empty input} است. در حلقه بعدی، عرض هر ستون بر اساس حداکثر طول مقدار در آن ستون محاسبه می‌شود؛ این کار با ترکیب \lr{\texttt{max}} و \lr{\texttt{generator expression}} انجام می‌شود که یک الگوی پایتونیک برای محاسبه حداکثر در یک مجموعه است. سپس خط عنوان، خط جداکننده (با \lr{-}‌ها) و سطرهای داده به‌ترتیب به خطوط اضافه می‌شوند. استفاده از \lr{\texttt{ljust(w)}} تضمین می‌کند که تمام مقادیر در یک ستون، عرض یکسان داشته باشند و جدول هم‌تراز باشد.

\subsection{تابع \_short\_identity و نمایش فشرده هویت}

تابع \lr{\texttt{\_short\_identity}} یک \lr{dict} هویت پالس را به یک رشته فشرده به فرمت \lr{bit/basis/bob/state} تبدیل می‌کند. این نمایش، خوانایی جدول ممیزی را به‌طور قابل‌توجهی افزایش می‌دهد؛ زیرا یک رشته کوتاه به‌جای چهار فیلد جداگانه نمایش داده می‌شود.

\begin{LTR}
	\begin{lstlisting}[language=Python]
		def _short_identity(d: Dict[str, Any]) -> str:
		"""Compact representation of an identity dict: 'bit/basis/bob/state'."""
		return (f"{d['alice_bit']}/{d['alice_basis']}/"
		f"{d['bob_basis']}/{d['state_label']}")
	\end{lstlisting}
\end{LTR}

این تابع با استفاده از \lr{f-string} و \lr{/} به‌عنوان جداکننده، یک رشته به فرمت \lr{0/0/1/|+⟩} تولید می‌کند که در آن اولین مقدار \lr{alice\_bit}، دومین \lr{alice\_basis}، سومین \lr{bob\_basis} و چهارمین \lr{state\_label} است. این فرمت، مشابه نمایش \lr{CSV} است اما با جداکننده \lr{/} که در خروجی متنی خواناتر است.

\subsection{تابع \_identity\_match و منطق مقایسه چهارمرحله‌ای}

تابع \lr{\texttt{\_identity\_match}} چهار \lr{dict} هویتی (\lr{s1}، \lr{s2}، \lr{s3}، \lr{s4}) را می‌گیرد و در صورتی که همه فیلدهای \lr{IDENTITY\_FIELDS} در هر چهار مرحله برابر باشند، مقدار \lr{"OK"} و در غیر این صورت، \lr{"FAIL"} برمی‌گرداند.

\begin{LTR}
	\begin{lstlisting}[language=Python]
		def _identity_match(s1, s2, s3, s4) -> str:
		"""Return 'OK' if all four stage views agree on every identity field."""
		for f in IDENTITY_FIELDS:
		if not (s1[f] == s2[f] == s3[f] == s4[f]):
		return "FAIL"
		return "OK"
	\end{lstlisting}
\end{LTR}

در این پیاده‌سازی، عبارت \lr{\texttt{s1[f] == s2[f] == s3[f] == s4[f]}} یک الگوی پایتونیک برای مقایسه زنجیره‌ای است که در آن، همه چهار مقدار باید برابر باشند. این عبارت معادل \lr{\texttt{s1[f] == s2[f] and s2[f] == s3[f] and s3[f] == s4[f]}} است اما خواناتر و کوتاه‌تر است. در صورت عدم تطابق در هر فیلد، تابع بلافاصله \lr{"FAIL"} برمی‌گرداند که الگوی \lr{early exit} است و از پردازش غیرضروری جلوگیری می‌کند.

\subsection{تابع \_measurement\_match و اعتبارسنجی پیش‌بینی}

تابع \lr{\texttt{\_measurement\_match}} دو پارامتر \lr{s1} (هویت پیش از \lr{dispatch}) و \lr{s3} (خروجی کارگر) را می‌گیرد و در صورتی که نتیجه اندازه‌گیری گزارش‌شده توسط کارگر با پیش‌بینی بر اساس \lr{s1} هم‌خوان باشد، \lr{"OK"} برمی‌گرداند.

\begin{LTR}
	\begin{lstlisting}[language=Python]
		def _measurement_match(s1, s3) -> str:
		"""Return 'OK' if the worker's measurement matches the S1 prediction."""
		exp = _expected_measurement(s1)
		if (s3["basis_match"] == exp["basis_match"] and
		s3["measured_bit"] == exp["measured_bit"] and
		s3["detected"] == exp["detected"]):
		return "OK"
		return "FAIL"
	\end{lstlisting}
\end{LTR}

در این تابع، ابتدا \lr{\texttt{\_expected\_measurement(s1)}} فراخوانی می‌شود تا پیش‌بینی نتیجه اندازه‌گیری بر اساس هویت پالس به دست آید. سپس سه فیلد کلیدی — \lr{\texttt{basis\_match}}، \lr{\texttt{measured\_bit}} و \lr{\texttt{detected}} — بین گزارش کارگر و پیش‌بینی مقایسه می‌شوند. این اعتبارسنجی، خطاهای احتمالی در منطق اندازه‌گیری کارگر را آشکار می‌سازد. به‌عنوان مثال، اگر کارگر به‌اشتباه \lr{\texttt{basis\_match = True}} را در حالی که پایه‌ها متفاوت هستند گزارش دهد، این تابع خطا را شناسایی می‌کند.

\subsection{تابع run\_trace و الگوی Closure برای خروجی}

تابع \lr{\texttt{run\_trace}} اصلی‌ترین تابع ماژول است. این تابع یک \lr{closure} به نام \lr{\texttt{out}} تعریف می‌کند که رشته‌ها را به لیست \lr{\texttt{lines}} اضافه می‌کند:

\begin{LTR}
	\begin{lstlisting}[language=Python]
		def run_trace(args) -> str:
		"""Run the 100-pulse trace and return the full text output."""
		lines: List[str] = []
		def out(s: str = ""):
		lines.append(s)
	\end{lstlisting}
\end{LTR}

این الگو چندین مزیت دارد: نخست، تمام خروجی در یک لیست جمع‌آوری می‌شود و در پایان به یک رشته با \lr{\texttt{"\textbackslash n".join(lines)}} تبدیل می‌گردد. این کار، از هزینه \lr{string concatenation} مکرر (که در پایتون از مرتبه $O(n^2)$ است) جلوگیری می‌کند. دوم، تابع \lr{\texttt{out}} به‌عنوان یک \lr{closure} به \lr{\texttt{lines}} دسترسی دارد و نیازی به ارسال \lr{\texttt{lines}} به عنوان پارامتر نیست. سوم، این الگو امکان تغییر آینده به \lr{file-like object} یا \lr{stream} را فراهم می‌سازد.

\subsection{تنظیم مسیر و الگوی Path Bootstrap}

در ابتدای ماژول، الگوی \lr{path bootstrap} پیاده‌سازی شده است:

\begin{LTR}
	\begin{lstlisting}[language=Python]
		SCRIPT_DIR = os.path.dirname(os.path.abspath(__file__))
		PROJECT_ROOT = os.path.dirname(SCRIPT_DIR)
		for candidate in (
		PROJECT_ROOT,
		os.path.join(PROJECT_ROOT, "src"),
		os.path.join(PROJECT_ROOT, "tests"),
		):
		if candidate not in sys.path:
		sys.path.insert(0, candidate)
	\end{lstlisting}
\end{LTR}

این الگو، تضمین می‌کند که ماژول بدون توجه به محل اجرای آن (از \lr{project root} یا از مسیر دیگر)، بتواند ماژول‌های \lr{\texttt{tests}} و \lr{\texttt{src}} را وارد کند. \lr{\texttt{SCRIPT\_DIR}} مسیر دایرکتوری حاوی اسکریپت است و \lr{\texttt{PROJECT\_ROOT}} یک سطح بالاتر است. سپس سه مسیر (ریشه پروژه، \lr{src}، \lr{tests}) به \lr{\texttt{sys.path}} اضافه می‌شوند. این الگو در اسکریپت‌هایی که به‌صورت مستقل اجرا می‌شوند و نیاز به واردسازی ماژول‌های پروژه دارند، استاندارد است.

\section{پیاده‌سازی ریاضی و الگوریتمی}

\subsection{الگوریتم انتخاب تصادفی پالس‌های ردیابی‌شده}

انتخاب ۱۰۰ پالس از جمعیت ۱۰۰۰ پالس، با بذر تصادفی مشخص، توسط تابع \lr{\texttt{\_select\_tracked\_pids}} انجام می‌شود. این تابع از \lr{\texttt{random.Random(seed)}} برای تولید یک \lr{PRNG} قابل بازتولید استفاده می‌کند و سپس با \lr{\texttt{random.sample}}، $k$ شناسه از جمعیت $N$ انتخاب می‌نماید. این الگوریتم، الگوریتم \lr{Fisher-Yates partial shuffle} است که در آن، فقط $k$ جایگشت اول محاسبه می‌شود:

\begin{equation}
	P(\text{pid } i \text{ selected}) = \frac{k}{N}, \quad \forall i \in \{0, 1, \ldots, N-1\}
\end{equation}

این الگوریتم تضمین می‌نماید که هر پالس با احتمال یکسان $k/N$ انتخاب شود و انتخاب‌ها مستقل باشند. با بذر مشخص، انتخاب قابل بازتولید است که برای آزمون‌های \lr{regression} حیاتی است.

\subsection{الگوریتم محاسبه Pairwise Inversions}

محاسبه تعداد \lr{pairwise inversions} در لیست ورودی، با الگوریتم ساده $O(n^2)$ انجام می‌شود:

\begin{LTR}
	\begin{lstlisting}[language=Python]
		n = len(tracked_arrival_order)
		inversions = 0
		for i in range(n):
		ai = tracked_arrival_order[i]
		for j in range(i + 1, n):
		if ai > tracked_arrival_order[j]:
		inversions += 1
		max_inv = n * (n - 1) // 2
	\end{lstlisting}
\end{LTR}

در این الگوریتم، برای هر زوج (\lr{i, j}) با $i < j$، بررسی می‌شود که آیا \lr{ai > arrival[j]} است. در صورت مثبت بودن، یک وارونگی شمرده می‌شود. حداکثر تعداد وارونگی‌ها برابر با $\binom{n}{2} = \frac{n(n-1)}{2}$ است که با عملگر \lr{//} (تقسیم صحیح) محاسبه می‌شود. این الگوریتم برای $n = 100$ کافی است ($10^4$ عملیات)، اما برای $n$ بزرگ‌تر، الگوریتم \lr{merge-sort-based} با پیچیدگی $O(n \log n)$ ترجیح داده می‌شود.

\subsection{الگوریتم یافتن Concrete Swap Examples}

برای نمایش شهودی \lr{reordering}، الگوریتمی برای یافتن نمونه‌های ملموس جابجایی پیاده‌سازی شده است:

\begin{LTR}
	\begin{lstlisting}[language=Python]
		swap_examples = []
		for i in range(len(tracked_arrival_order) - 1):
		a = tracked_arrival_order[i]
		b = tracked_arrival_order[i + 1]
		if a > b:  # inversion: larger pid arrived before smaller
		swap_examples.append((i, i + 1, a, b))
		if len(swap_examples) >= 5:
		break
	\end{lstlisting}
\end{LTR}

این الگوریتم، زوج‌های مجاور (\lr{adjacent pairs}) را بررسی می‌کند و در صورتی که \lr{pid} بزرگ‌تر قبل از \lr{pid} کوچک‌تر آمده باشد، آن را به‌عنوان یک جابجایی ثبت می‌نماید. این الگوریتم، فقط \lr{adjacent inversions} را شناسایی می‌کند که زیرمجموعه‌ای از \lr{pairwise inversions} هستند. هدف از این الگوریتم، ارائه نمونه‌های ملموس و قابل درک برای کاربر است؛ زیرا مشاهده دو \lr{pid} که جابجا شده‌اند، شهودی‌تر از عدد انتزاعی \lr{inversion count} است. حداکثر ۵ نمونه نمایش داده می‌شود تا خروجی بیش از حد طولانی نشود.

\subsection{الگوریتم محاسبه Position Shifts}

برای تحلیل کمی جابجایی هر پالس، الگوریتم محاسبه \lr{position shifts} پیاده‌سازی شده است:

\begin{LTR}
	\begin{lstlisting}[language=Python]
		sorted_idx = {pid: i for i, pid in enumerate(tracked_pids)}
		arrived_idx = {pid: i for i, pid in enumerate(tracked_arrival_order)}
		shifts = [(pid, sorted_idx[pid], arrived_idx[pid], arrived_idx[pid] - sorted_idx[pid])
		for pid in tracked_pids]
		max_up = max(shifts, key=lambda x: x[3])
		max_down = min(shifts, key=lambda x: x[3])
	\end{lstlisting}
\end{LTR}

در این الگوریتم، ابتدا دو دیکشنری \lr{\texttt{sorted\_idx}} و \lr{\texttt{arrived\_idx}} ساخته می‌شوند که برای هر \lr{pid}، موقعیت آن در لیست مرتب (ورودی) و در لیست ورودی (خروجی) را نگه می‌دارند. سپس برای هر پالس، اختلاف این دو موقعیت (\lr{shift}) محاسبه می‌شود. \lr{shift} مثبت به این معناست که پالس در خروجی زودتر از ورودی آمده (حرکت به بالا) و \lr{shift} منفی به این معناست که پالس دیرتر آمده (حرکت به پایین). \lr{\texttt{max\_up}} و \lr{\texttt{max\_down}} به‌ترتیب بیشترین حرکت به بالا و بیشترین حرکت به پایین را نشان می‌دهند. این الگوریتم، الگوی \lr{dictionary comprehension} و \lr{list comprehension} را به‌خوبی به نمایش می‌گذارد.

\subsection{الگوریتم Stage-Level Mismatch Counting}

برای محاسبه آمار ناهماهنگی در هر مرحله، الگوریتم زیر پیاده‌سازی شده است:

\begin{LTR}
	\begin{lstlisting}[language=Python]
		for pid in tracked_pids:
		if pid not in receipts_by_pid:
		missing_receipts += 1
		continue
		if pid not in rows_by_pid:
		missing_rows += 1
		continue
		s1 = snapshot[pid]
		s2 = receipts_by_pid[pid]["received_item"]
		s3 = rows_by_pid[pid]
		s4 = rows_by_pid[pid]
		if any(s1[f] != s2[f] for f in IDENTITY_FIELDS):
		s1_s2_mismatch += 1
		if any(s2[f] != s3[f] for f in IDENTITY_FIELDS):
		s2_s3_mismatch += 1
		if any(s3[f] != s4[f] for f in IDENTITY_FIELDS):
		s3_s4_mismatch += 1
		exp = _expected_measurement(s1)
		if (s3["basis_match"] != exp["basis_match"] or
		s3["measured_bit"] != exp["measured_bit"] or
		s3["detected"] != exp["detected"]):
		measurement_mismatch += 1
	\end{lstlisting}
\end{LTR}

در این الگوریتم، برای هر \lr{pid} در \lr{tracked\_pids}، ابتدا وجود \lr{receipt} و \lr{row} بررسی می‌شود. در صورت عدم وجود، شمارنده \lr{missing} افزایش می‌یابد و با \lr{\texttt{continue}} به پالس بعدی می‌رود. در غیر این صورت، چهار \lr{dict} هویتی \lr{s1}، \lr{s2}، \lr{s3}، \lr{s4} استخراج می‌شوند و برای هر جفت متوالی (\lr{S1-S2}، \lr{S2-S3}، \lr{S3-S4})، تطابق فیلدهای هویتی بررسی می‌گردد. تابع \lr{\texttt{any}} با \lr{generator expression}، الگوی پایتونیک برای بررسی وجود عدم تطابق است. در پایان، تطابق نتیجه اندازه‌گیری نیز بررسی می‌شود. این الگوریتم، الگوی \lr{fail-fast} را رعایت می‌کند: در صورت فقدان داده‌های پایه‌ای (\lr{receipt} یا \lr{row})، بقیه بررسی‌ها برای آن پالس انجام نمی‌شود.

\subsection{الگوریتم Verdict و تفسیر نهایی}

در پایان، الگوریتم \lr{verdict} نتیجه نهایی را اعلام می‌کند:

\begin{LTR}
	\begin{lstlisting}[language=Python]
		total_mismatch = (missing_receipts + missing_rows + s1_s2_mismatch +
		s2_s3_mismatch + s3_s4_mismatch + measurement_mismatch)
		if total_mismatch == 0:
		out("  VERDICT:  All 100 tracked pulses preserved identity through all 4 stages.")
		out(f"           {inversions} pairwise inversions occurred (reordering), yet NOT A")
		out("           SINGLE pulse's identity was corrupted. Identity is bound to")
		out("           pulse_id, NOT to row position -- exactly as the contract requires.")
		else:
		out(f"  VERDICT:  {total_mismatch} corruption(s) detected across the 4 stages.")
		out("           See the per-pulse audit table above for details.")
	\end{lstlisting}
\end{LTR}

این الگوریتم، مجموع تمام ناهماهنگی‌ها را محاسبه می‌کند و در صورت صفر بودن، پیام موفقیت با تأکید بر تمایز بین \lr{reordering} (که قابل‌قبول است) و \lr{corruption} (که باگ است) چاپ می‌نماید. این تمایز، پیام کلیدی ماژول است: \lr{reordering} یک ویژگی ذاتی شبیه‌سازی موازی است و در صورتی که هویت به \lr{pulse\_id} بسته شده باشد، هیچ مشکلی ایجاد نمی‌کند.

\subsection{الگوریتم CLI با argparse}

تابع \lr{\texttt{main}} از \lr{\texttt{argparse}} برای پیکربندی پارامترها استفاده می‌کند:

\begin{LTR}
	\begin{lstlisting}[language=Python]
		def main():
		parser = argparse.ArgumentParser(
		description="100-random-pulse end-to-end state trace",
		formatter_class=argparse.RawDescriptionHelpFormatter,
		)
		parser.add_argument("--workers", type=int, default=4)
		parser.add_argument("--population", type=int, default=1000)
		parser.add_argument("--tracked", type=int, default=100)
		parser.add_argument("--seed", type=int, default=42)
		parser.add_argument("--first", type=int, default=100)
		parser.add_argument("--output", type=str, default=None)
		args = parser.parse_args()
		
		text = run_trace(args)
		print(text)
		
		default_path = os.path.join(PROJECT_ROOT, "download",
		"pulse_tracking_trace.txt")
		os.makedirs(os.path.dirname(default_path), exist_ok=True)
		with open(default_path, "w") as f:
		f.write(text)
		print(f"\n[trace also written to {default_path}]", file=sys.stderr)
	\end{lstlisting}
\end{LTR}

این الگوریتم، الگوی استاندارد \lr{CLI} در پایتون است. هر پارامتر با \lr{\texttt{add\_argument}} تعریف می‌شود و دارای مقدار پیش‌فرض است. \lr{\texttt{RawDescriptionHelpFormatter}} امکان قالب‌بندی متن راهنما را فراهم می‌سازد. پس از اجرای \lr{\texttt{run\_trace}}، خروجی به \lr{stdout} چاپ می‌شود و همچنین به فایل پیش‌فرض در \lr{\texttt{download/}} نوشته می‌گردد. در صورت ارائه \lr{\texttt{-{}-output}}، خروجی به فایل اضافی نیز نوشته می‌شود. این الگو، تضمین می‌کند که خروجی همیشه برای بازبینی بعدی ذخیره می‌شود، حتی اگر کاربر فایل خاصی مشخص نکند.

\section{ارسال و دریافت با سایر ماژول‌ها}

\subsection{وابستگی به tests/test\_main\_optimized\_pulse\_tracking}

اصلی‌ترین وابستگی این ماژول، به \lr{\texttt{tests/test\_main\_optimized\_pulse\_tracking.py}} است. این وابستگی شامل پنج مورد است: ثابت \lr{\texttt{IDENTITY\_FIELDS}} و چهار تابع \lr{\texttt{\_make\_pulse\_population}}، \lr{\texttt{\_select\_tracked\_pids}}، \lr{\texttt{\_run\_tracking\_dispatch}} و \lr{\texttt{\_expected\_measurement}} و \lr{\texttt{\_tracking\_slow\_worker}}. این الگوی \lr{reuse-by-import}، تضمین می‌کند که اسکریپت ممیزی و آزمون واحد، از همان ماشین‌آلات استفاده کنند و نتایج آن‌ها همیشه هم‌خوان باشند. در صورت تغییر در یکی از این توابع (مثلاً افزودن فیلد جدید به \lr{\texttt{IDENTITY\_FIELDS}})، هر دو به‌طور خودکار به‌روز می‌شوند.

این وابستگی، یک وابستگی \lr{tight coupling} است: تغییر در ساختار \lr{dict} پالس یا در امضای توابع، می‌تواند منجر به شکست این ماژول شود. برای کاهش این \lr{coupling}، می‌توان توابع کمکی را به یک ماژول جداگانه (مانند \lr{\texttt{pulse\_tracking\_utils.py}}) منتقل کرد که هم اسکریپت ممیزی و هم آزمون واحد از آن استفاده کنند. این الگو، الگوی \lr{shared utility module} است که در پروژه‌های بالغ به‌کررات به‌کار می‌رود.

\subsection{ورودی از main\_optimized\_dispatch}

این ماژول تابع \lr{\texttt{dispatch\_work\_items}} را از \lr{\texttt{main\_optimized\_dispatch}} وارد می‌کند. این تابع، موتور اصلی \lr{dispatch} فریم‌ورک است که کارهای شبیه‌سازی را بین چند کارگر توزیع می‌کند. وابستگی به این تابع، تضمین می‌کند که اسکریپت ممیزی، همان مسیر اجرای واقعی فریم‌ورک را آزمون می‌کند و نه یک مسیر آزمایشگاهی مصنوعی.

\subsection{خروجی به stdout و فایل}

خروجی این ماژول به دو شکل است: نخست، چاپ به \lr{stdout} با \lr{\texttt{print(text)}} که امکان مشاهده زنده را فراهم می‌سازد. دوم، نوشتن به فایل پیش‌فرض در مسیر \lr{\texttt{/home/z/my-project/download/pulse\_tracking\_trace.txt}} با \lr{\texttt{open(default\_path, "w")}}. در صورت ارائه \lr{\texttt{-{}-output}}، خروجی به فایل اضافی نیز نوشته می‌شود. این الگو، تضمین می‌کند که خروجی همیشه برای بازبینی بعدی ذخیره می‌شود، حتی اگر کاربر آن را به فایل خاصی هدایت نکرده باشد.

\subsection{ارتباط با فریم‌ورک آزمون}

این ماژول و ماژول آزمون \lr{\texttt{tests/test\_main\_optimized\_pulse\_tracking.py}} دو لایه مکمل اعتبارسنجی را تشکیل می‌دهند. لایه آزمون، اعتبارسنجی خودکار را با \lr{pytest} فراهم می‌سازد و در \lr{CI/CD pipeline} اجرا می‌شود. لایه اسکریپت ممیزی (این ماژول)، خروجی قابل مشاهده و قابل تحلیل توسط انسان را فراهم می‌کند و برای \lr{debug} و تحلیل دستی استفاده می‌شود. این دو لایه، نقش‌های مکمل دارند: لایه آزمون، \lr{gate} خودکار را برای \lr{merge} فراهم می‌سازد و لایه اسکریپت، امکان تحلیل عمیق در صورت بروز مشکل را فراهم می‌نماید.

\subsection{خروجی قابل تحلیل توسط ابزارهای خارجی}

خروجی متنی این ماژول، به‌گونه‌ای طراحی شده که توسط ابزارهای تحلیل متنی (مانند \lr{grep}، \lr{awk} یا اسکریپت‌های \lr{Python}) قابل پردازش باشد. جدول‌ها با \lr{spaces} هم‌تراز شده‌اند و جداکننده‌ها واضح هستند. در صورت نیاز، می‌توان خروجی را به فرمت \lr{JSON} یا \lr{CSV} نیز تبدیل کرد تا توسط ابزارهای تحلیل داده (مانند \lr{pandas}) پردازش شود. این انعطاف‌پذیری، در پروژه‌های تحقیقاتی که نیاز به تحلیل آماری نتایج شبیه‌سازی دارند، ارزشمند است.

\subsection{جمع‌بندی ارتباطات}

به‌طور خلاصه، این ماژول در نقش یک \lr{forensic tracer} مستقل عمل می‌کند که از طریق سه کانال با سایر بخش‌های فریم‌ورک در ارتباط است: ورودی توابع کمکی از \lr{\texttt{tests/test\_main\_optimized\_pulse\_tracking.py}}، ورودی موتور \lr{dispatch} از \lr{\texttt{main\_optimized\_dispatch}}، و خروجی به \lr{stdout} و فایل \lr{\texttt{pulse\_tracking\_trace.txt}}. این معماری، جداسازی واضحی بین لایه ممیزی و موتور شبیه‌سازی برقرار می‌سازد و امکان توسعه مستقل هر یک را فراهم می‌نماید. خروجی این ماژول، نه‌تنها برای توسعه‌دهندگان (در فرآیند \lr{debug})، بلکه برای بازبینان کریپتوگرافیک (در فرآیند \lr{security audit}) نیز مفید است، زیرا یک \lr{audit trail} قابل ممیزی از حفظ هویت پالس‌ها در طول خط لوله شبیه‌سازی فراهم می‌سازد.
